\documentclass[11pt,a4paper]{article}

\usepackage[spanish,es-noquoting]{babel}
\usepackage[utf8]{inputenc}
\usepackage[T1]{fontenc}
\usepackage{lmodern}
\usepackage{geometry}
\usepackage{booktabs}
\usepackage{tabularx}
\usepackage{longtable}
\usepackage{array}
\usepackage{float}
\usepackage{enumitem}
\usepackage{hyperref}

\geometry{margin=2.5cm}
\setlength{\parskip}{0.6em}
\setlength{\parindent}{0pt}
\hypersetup{colorlinks=true, linkcolor=black, urlcolor=blue}

\title{Proyecto 2 - Data Mining\\\large Informe Final Integrado (Semanas 1--4)}
\author{Javier Espana \and Angel Esquit \and Roberto Barreda \and Diego Lopez}
\date{24 abril 2026}

\begin{document}

\maketitle

\begin{center}
Universidad del Valle de Guatemala\\
Curso: Data Mining\\
Dataset: Wisconsin Breast Cancer Dataset\\
Problema: Clasificacion binaria de tumor mamario (2=Benigno, 4=Maligno)
\end{center}

\hrule

\section{Resumen ejecutivo}

Este proyecto desarrollo y comparo modelos de clasificacion para apoyo diagnostico de cancer de mama a partir de rasgos citologicos del Wisconsin Breast Cancer Dataset. El flujo completo se ejecuto en cuatro semanas: comprension y EDA, preparacion/modelos base, validacion y tuning, e interpretabilidad/errores.

El resultado final fue la seleccion de \textbf{SVM optimizado con kernel RBF en Escenario A (todas las variables)}, debido a su desempeno clinicamente prioritario: \textbf{Recall de clase maligna = 1.0000} y \textbf{FN = 0} en test. La decision se tomo con trazabilidad completa y comparacion formal contra LR y KNN optimizados, aceptando el trade-off de mayor FP frente a la reduccion total de FN.

\section{Introduccion y planteamiento del problema}

La deteccion temprana de lesiones malignas en mama es un problema critico por su impacto en pronostico y tratamiento. En este contexto, el principal riesgo operacional de un modelo es omitir un caso maligno (falso negativo).

Por ello, el proyecto se diseno con una prioridad clara:

\begin{itemize}[leftmargin=1.2cm]
\item minimizar falsos negativos de la clase maligna,
\item sin descuidar estabilidad general de desempeno y consistencia estadistica.
\end{itemize}

\section{Datos y descripcion del dataset}

Se utilizo \texttt{Datos.csv} con 683 filas originales y 11 columnas (9 rasgos citologicos, identificador y target). La variable objetivo es \texttt{Class}, con distribucion original:

\begin{itemize}[leftmargin=1.2cm]
\item Benigno (2): 444 casos (65.01\%)
\item Maligno (4): 239 casos (34.99\%)
\end{itemize}

Existe desbalance moderado (1.86:1), por lo que accuracy no fue usada como criterio unico de decision.

\section{Semana 1: Comprension del problema y EDA}

\subsection{Hallazgos principales}

\begin{enumerate}[leftmargin=1.2cm]
\item \textbf{Separabilidad}: se observaron distribuciones diferenciadas entre benignos y malignos, especialmente en variables de uniformidad y nucleos.
\item \textbf{Variables dominantes}: \texttt{Bare Nuclei}, \texttt{Uniformity of Cell Size} y \texttt{Uniformity of Cell Shape} mostraron mayor asociacion con la clase.
\item \textbf{Colinealidad}: alta correlacion entre \texttt{Uniformity of Cell Size} y \texttt{Uniformity of Cell Shape} (r=0.91).
\item \textbf{Outliers clinicos}: se detectaron valores extremos, sobre todo en casos malignos, que se mantuvieron por su posible valor diagnostico.
\end{enumerate}

\subsection{Hipotesis de trabajo (H1--H5)}

Se formularon hipotesis sobre relevancia de variables, utilidad de modelos no lineales, prioridad de recall maligna, impacto de colinealidad y conservacion de outliers. Estas hipotesis guiaron las decisiones de Semanas 2 y 3.

\section{Semana 2: Preparacion de datos y modelos base}

\subsection{Limpieza y version de modelado}

Pasos aplicados:

\begin{itemize}[leftmargin=1.2cm]
\item conversion y manejo de invalidez/nulos (incluyendo \texttt{?} en \texttt{Bare Nuclei}),
\item eliminacion de 8 duplicados exactos,
\item exclusion de \texttt{Sample code number} para modelado.
\end{itemize}

Resultado:

\begin{itemize}[leftmargin=1.2cm]
\item \texttt{Datos\_limpio.csv}: 675 filas (con ID)
\item \texttt{Datos\_modelado.csv}: 675 filas (sin ID)
\end{itemize}

\subsection{Escenarios y particion}

\begin{itemize}[leftmargin=1.2cm]
\item Escenario A: 9 variables.
\item Escenario B: top 3 variables (respuesta metodologica al hallazgo de colinealidad y relevancia).
\item Split 80/20 estratificado con \texttt{random\_state=42}.
\end{itemize}

\subsection{Modelos base}

Se entrenaron sin tuning:

\begin{itemize}[leftmargin=1.2cm]
\item Regresion Logistica (L2)
\item KNN (k=5)
\item SVM (RBF)
\end{itemize}

El mejor baseline ya mostraba ventaja de SVM en sensibilidad de clase maligna.

\section{Semana 3: Validacion, optimizacion y seleccion final}

\subsection{Diseno experimental}

Se aplico protocolo unico para comparabilidad:

\begin{itemize}[leftmargin=1.2cm]
\item \texttt{StratifiedKFold} de 5 folds,\
\item pipelines con escalado interno,\
\item test set reservado para verificacion final,\
\item metrica objetivo de tuning: recall de clase maligna (Class=4).
\end{itemize}

\subsection{Resultados optimizados (resumen oficial)}

\begin{table}[H]
\centering
\caption{Modelos optimizados (extracto de \texttt{19\_optimized\_only.csv})}
\begin{tabular}{llcccc}
\toprule
Escenario & Modelo & Recall CV & Recall Test & Precision Test & FN Test \\
\midrule
A & SVM & 1.0000 & 1.0000 & 0.9038 & 0 \\
A & LR & 0.9789 & 0.9574 & 0.9375 & 2 \\
A & KNN & 0.9632 & 0.9149 & 0.9348 & 4 \\
B & SVM & 0.9735 & 0.9787 & 0.9388 & 1 \\
B & LR & 0.9471 & 0.9149 & 0.9348 & 4 \\
B & KNN & 0.9578 & 0.9149 & 0.9348 & 4 \\
\bottomrule
\end{tabular}
\end{table}

\subsection{Seleccion final}

\textbf{Modelo seleccionado: SVM optimizado, Escenario A.}

Justificacion:

\begin{enumerate}[leftmargin=1.2cm]
\item Recall maligna maximo (1.0000).
\item FN minimo absoluto (0).
\item Consistencia CV-test adecuada.
\item Coherencia con objetivo clinico del proyecto.
\end{enumerate}

Trade-off explicitado:

\begin{itemize}[leftmargin=1.2cm]
\item Frente a LR-A optimizada, SVM-A evita 2 FN adicionales.
\item A cambio incrementa 2 FP adicionales.
\item En criterio clinico definido, este intercambio es aceptable.
\end{itemize}

\section{Semana 4: Interpretabilidad y analisis de errores}

\subsection{Interpretabilidad del modelo final}

Como SVM-RBF no ofrece importancias directas, se implemento interpretacion complementaria:

\begin{itemize}[leftmargin=1.2cm]
\item ranking por Random Forest (importancia por impureza),
\item ranking por Regresion Logistica (magnitud absoluta de coeficientes),
\item ranking promedio para robustez de lectura.
\end{itemize}

Variables dominantes recurrentes:

\begin{itemize}[leftmargin=1.2cm]
\item \texttt{Bare Nuclei}
\item \texttt{Uniformity of Cell Shape}
\item \texttt{Uniformity of Cell Size}
\end{itemize}

Esto confirma coherencia longitudinal con EDA y seleccion de escenarios.

\subsection{Analisis de errores del modelo final}

Reproduciendo SVM-A con hiperparametros oficiales de Semana 3:

\begin{itemize}[leftmargin=1.2cm]
\item TN=83, FP=5, FN=0, TP=47
\item Error total en test: 5 casos, todos FP
\end{itemize}

No se observaron FN, por lo que el analisis se concentro en FP y sus patrones. Los FP muestran elevaciones en variables como \texttt{Bare Nuclei}, \texttt{Single Epithelial Cell Size}, \texttt{Normal Nucleoli}, \texttt{Uniformity of Cell Size} y \texttt{Marginal Adhesion}, sugiriendo casos benignos cercanos a la frontera de malignidad.

\subsection{Implicacion}

El comportamiento final del modelo es consistente con la politica de riesgo adoptada: sensibilidad maxima con incremento moderado de revisiones adicionales por FP.

\section{Discusion integrada del proyecto}

\subsection{Coherencia metodologica}

La narrativa tecnica se mantuvo consistente:

\begin{itemize}[leftmargin=1.2cm]
\item Semana 1 define hallazgos e hipotesis,
\item Semana 2 implementa escenarios y baselines,
\item Semana 3 valida, optimiza y selecciona,
\item Semana 4 explica decisiones del modelo final y limita riesgos.
\end{itemize}

\subsection{Fortalezas}

\begin{enumerate}[leftmargin=1.2cm]
\item Control experimental homogeno y reproducible.
\item Priorizacion explicita de metrica clinicamente relevante.
\item Justificacion cuantitativa de trade-offs.
\item Cierre de trazabilidad entre EDA, tuning y seleccion final.
\end{enumerate}

\subsection{Limitaciones}

\begin{enumerate}[leftmargin=1.2cm]
\item Dataset moderado en tamano, sin validacion externa.
\item SVM-RBF requiere interpretabilidad indirecta.
\item Persisten FP en frontera de decision.
\end{enumerate}

\section{Conclusiones finales}

\begin{enumerate}[leftmargin=1.2cm]
\item El objetivo principal del proyecto se cumplio: construir un clasificador con minima omision de malignos.
\item SVM optimizado en Escenario A fue el unico candidato con FN=0 en test bajo protocolo comparable.
\item El modelo final no se eligio por accuracy aislada, sino por criterio clinico y evidencia multi-metrica.
\item Las variables mas relevantes del sistema final son coherentes con los hallazgos citologicos iniciales.
\item El cierre de Semana 4 fortalece la defendibilidad tecnica del proyecto completo.
\end{enumerate}

\section{Recomendaciones finales}

\begin{enumerate}[leftmargin=1.2cm]
\item Mantener SVM-A como modelo base si la prioridad sigue siendo minimizar FN.
\item Acompanhar despliegue con protocolo de segunda revision para positivos de baja confianza.
\item Recalibrar umbral si el contexto requiere reducir FP sin perder sensibilidad critica.
\item Realizar validacion externa en cohortes adicionales antes de uso operativo fuera del curso.
\item Mantener reporte de metricas en bloque (Recall maligna, FN, Precision, AUC, FP), evitando decisiones por una sola metrica.
\end{enumerate}

\section{Entregables consolidados}

\begin{itemize}[leftmargin=1.2cm]
\item Informes semanales:\
\texttt{docs/01\_Semana\_01\_Informe.tex}\
\texttt{docs/02\_Semana\_02\_Informe.tex}\
\texttt{docs/03\_Semana\_03\_Informe.tex}\
\texttt{docs/04\_Semana\_04\_Informe.tex}
\item Informe final integrado: \texttt{docs/06\_Informe\_Final\_Proyecto.tex}
\item Notebooks tecnicos: carpeta \texttt{notebook/}
\item Resultados CSV de soporte: \texttt{data/resultados\_semana3\_csv/} y \texttt{data/resultados\_semana4\_csv/}
\end{itemize}

\end{document}
